\section{بار کاری پردازنده (\lr{cpu})}
\fig{fig/cpu.png}{تأخیر میانگین و \lr{P99} در آزمون \lr{cpu}}
\begin{center}
\begin{tabular}{lrrrr}
\toprule
محیط & \lr{avg (ms)} & انحراف معیار & \lr{P99 (ms)} & انحراف معیار \\
\midrule
\lr{Bare Metal (WSL)} & \lr{2565.28} & \lr{73.24} & \lr{2729.99} & \lr{98.30} \\
\lr{VM} & \lr{2775.21} & \lr{288.35} & \lr{3110.37} & \lr{576.40} \\
\lr{Docker} & \lr{2533.83} & \lr{58.71} & \lr{2665.33} & \lr{100.69} \\
\bottomrule
\end{tabular}
\end{center}
\textbf{تحلیل.} \lr{Docker} و \lr{Bare Metal} عملاً هم‌سطح‌اند (اختلاف میانگین حدود ۱٪ و کمتر از انحراف معیار تکرارها). دلیل آن است که کانتینر فقط یک فرآیند ایزوله‌شده (با \lr{namespace} و \lr{cgroup}) روی همان هستهٔ میزبان است و کد پردازشی مستقیماً روی \lr{CPU} اجرا می‌شود؛ هیچ لایهٔ ترجمه یا \lr{trap} اضافه‌ای وجود ندارد. در \lr{VM} میانگین حدود ۸٪ و \lr{P99} حدود ۱۴٪ نسبت به \lr{Bare Metal} بیشتر است، اما این اختلاف عمدتاً ناشی از تکرار سوم \lr{VM} (میانگین \lr{3107.97} و \lr{P99} برابر \lr{3773.42}) است. دو تکرار اول \lr{VM} با \lr{Bare Metal} تقریباً یکسان‌اند. این بالا رفتن ناگهانی با زمان‌بندی دو لایه (زمان‌بند میهمان و زمان‌بند \lr{hypervisor}) و رقابت \lr{vCPU} با سایر فعالیت‌های میزبان سازگار است. با سه تکرار نمی‌توان از نظر آماری نتیجه گرفت که سربار \lr{VM} ثابت و قطعی است؛ فقط می‌توان گفت \lr{VM} تغییرپذیری بیشتری دارد.
